\section{The weak rank condition}\label{sec:rank}

The size estimates are now available, but Proposition~\ref{prop:criterion}
also requires the weak rank condition \eqref{eq:weak-rank}. We prove it
first for the derivative rows $Q_k^{[p]}(-1)$, before removing the odd
columns. Section~\ref{sec:conclusion} then passes it to the integer forms
in even beta values.
The rows have zero entries for $i>b+h$: this holds for $C^{[p]}$ and
is preserved by \eqref{eq:Q}. We can therefore work with the system
truncated at component $b+h$.

The proof adapts \cite[Section~4.5]{F26}. A forbidden kernel vector would
give solutions whose polynomial combinations have many vanishing
derivatives at $-1$. We convert these into remainders for a single
polynomial row, prove that the remainders are independent, and apply
Proposition~\ref{prop:multiplicity} to contradict their total vanishing.
The inequality below is the one obtained from that final count.

\begin{proposition}\label{prop:rank}
Assume
\begin{equation}\label{eq:rank-condition}
 (h+1)(\kappa-2r)+\omega>a.
\end{equation}
For all sufficiently large admissible $n$, if
$x_0,\ldots,x_{b+h}\in\Q$ satisfy
\[
 \sum_{i=0}^{b+h}Q_{k,i}^{[p]}(-1)x_i=0
 \quad(0\leq p\leq h,\quad 2M+2\leq k\leq\kappa n),
\]
then $x_0=0$.
\end{proposition}

\subsection{Taylor corrections and a common polynomial row}

Suppose $x_0\neq0$, and scale so that $x_0=1$.
Let $g_0,\ldots,g_{b+h}$ solve \eqref{eq:base-system} near $-1$ with
$g_i(-1)=x_i$. In particular $g_0=1$.
The kernel equations say that certain high derivatives of
$\sum_i C_i^{[p]}g_i$ vanish at $-1$. To turn this into a zero of high
order, we must also remove the lower Taylor coefficients.
Choose $T_p\in\Q[z]$ of degree at most $2M$ to cancel the Taylor polynomial
through degree $2M$ of $\sum_i C_i^{[p]}g_i$ at $-1$.
The universal identity \eqref{eq:universal-row} and the assumed vanishing
of the higher derivatives give
\begin{equation}\label{eq:f-zero}
 f_p(z):=T_p(z)+\sum_{i=1}^{b+h}C_i^{[p]}(z)g_i(z)
          =O((z+1)^{\kappa n}).
\end{equation}
The derivatives of orders $0,\ldots,2M$ vanish by the choice of $T_p$;
those of orders $2M+1,\ldots,\kappa n-1$ vanish by hypothesis.

The functions $f_p$ still have different coefficient rows. The
multiplicity theorem instead needs one fixed polynomial row applied to
several solutions of the same system. Following Fischler, we arrange
this by combining the $f_p$ with powers of a local branch
$L(z)=-\log z$ near $-1$:
\[
 \rho_q(z)=\sum_{p=0}^q\binom qp L(z)^{q-p}f_p(z)
 \qquad(0\leq q\leq h).
\]
Each $\rho_q$ retains the zero of order at least $\kappa n$ at $-1$.
The common polynomial row will have $h+1$ entries for the corrections
$T_p$ and $b$ entries for the original polynomials $P_i$:
\begin{equation}\label{eq:S-common}
 S_i(z)=
 \begin{cases}
 T_{h-i}(z)/(h-i)!,&0\leq i\leq h,\\
 z^M P_{i-h}(z),&h+1\leq i\leq b+h,
 \end{cases}
\end{equation}
The dependence on $q$ is placed in the solution vectors. Define
\begin{equation}\label{eq:Yq}
 y_{i,q}(z)=
 \begin{cases}
 0,&i<h-q,\\[2pt]
 \displaystyle\frac{q!}{(i+q-h)!}L(z)^{i+q-h},&h-q\leq i\leq h,\\[6pt]
 \displaystyle\sum_{p=0}^q\binom qp L(z)^{q-p}
            (-1)^p(i-h)_p g_{i-h+p}(z),&i>h.
 \end{cases}
\end{equation}
Then $\rho_q=\sum_i S_i y_{i,q}$.
All vectors $Y_q=(y_{0,q},\ldots,y_{b+h,q})$ solve the same system
\begin{align}
 y_0'&=0,\notag\\
 y_i'&=-y_{i-1}/z &&(1\leq i\leq h),\notag\\
 y_{h+1}'&=\frac{2}{z(1-z)}y_h-\frac1{2z}y_{h+1},
                                                   \label{eq:enlarged-system}\\
 y_i'&=-y_{i-1}/z-y_i/(2z) &&(i\geq h+2).\notag
\end{align}
The first $h+1$ components are powers of $L$; the components with $i>h$
contain the original solution $g$.
On differentiating the latter components, terms from the derivatives of
$L^{q-p}$ combine with those from the derivatives of $g_{i-h+p}$.
After reindexing, their coefficients simplify by the identity
\[
 (j)_p-p(j)_{p-1}=(j-1)_p
 \qquad(j\geq2,\ p\geq1).
\]
The $p=0$ factor is $1$. This gives the equation at index $i=h+j$.
When $j=1$, the corresponding factors vanish for $p\geq1$; the remaining
source is $2L^q/[z(1-z)]$, giving the equation for $y_{h+1}$.
The diagonal term remains $-y_i/(2z)$ in every summand.
The equations for $i\leq h$ follow by differentiating powers of $L$.

\subsection{Quadratic change of variable}

The common diagonal term $-y_i/(2z)$ in the components $i>h$ can be
removed by multiplying those components by $\sqrt z$. We introduce
$w$ with $z=w^2$ so that this change is rational in the new variable
and the coefficient row remains polynomial. Set
\[
 v_i(w)=
 \begin{cases}
 y_i(w^2),&i\leq h,\\
 w\,y_i(w^2),&i>h.
 \end{cases}
\]
The diagonal terms cancel, leaving
\begin{align}
 v_0'&=0,\notag\\
 v_i'&=-2v_{i-1}/w &&(i\geq1,\ i\neq h+1),\label{eq:lifted-system}\\
 v_{h+1}'&=4v_h/(1-w^2).\notag
\end{align}
The common coefficient polynomials become
\begin{equation}\label{eq:lifted-polynomials}
 \widetilde S_i(w)=
 \begin{cases}
 S_i(w^2),&i\leq h,\\
 w^{2M-1}P_{i-h}(w^2),&i>h.
 \end{cases}
\end{equation}
Their degrees are at most $4M$: those with $i\leq h$ have degree at
most $4M$, and those with $i>h$ have degree at most
$2(M+n)-1\leq4M-1$.
For the solutions obtained from $Y_q$, the remainder
$\sum_i\widetilde S_i v_i$ equals $\rho_q(w^2)$.
It vanishes to order at least $\kappa n$ at each of $\ii$ and $-\ii$.

The change of variable increases the polynomial degrees, but it also
gives two points at which to count the high-order zeros.
The same polynomials \eqref{eq:lifted-polynomials} are used at both
points; only the local solution families differ.
This is why the Taylor corrections were constructed in $z$ before
making the change of variable.

\subsection{Independence by residues}

To count these zeros with Proposition~\ref{prop:multiplicity}, we must
know that no nonzero constant combination of the $h+1$ remainders
vanishes identically. The triangular form of the new system makes this
possible. Each derivative link has a nonzero residue, whereas the
derivative of a rational function has zero residue at every finite point.
The following lemma uses this to rule out rational-function relations
between successive components. It abstracts the argument in
\cite[Section~4.5, Step~4]{F26}. Different links of the system may use
different finite poles.

\begin{lemma}\label{lem:chain}
Suppose $v_0$ is a nonzero constant and
$v_j'=a_j(w)v_{j-1}$ for $1\leq j\leq t$, where each rational function
$a_j$ has a nonzero residue at some finite point.
Then $v_0,\ldots,v_t$ are linearly independent over $\C(w)$.
\end{lemma}

\begin{proof}
Clear the denominators of a putative rational relation and write it as
$\sum_{i=0}^{\alpha}U_i v_i=0$, with $U_\alpha\neq0$ and $\alpha\geq1$.
Let $D=\deg U_\alpha$. Differentiating $D+1$ times eliminates
$v_\alpha$. The coefficient of $v_{\alpha-1}$ is
\begin{equation}\label{eq:residue-coefficient}
 U_{\alpha-1}^{(D+1)}
   +\sum_{j=0}^{D}(a_\alpha U_\alpha^{(j)})^{(D-j)}.
\end{equation}
This follows by differentiating $U_\alpha v_\alpha$ successively and
retaining the two highest components.
All summands except $j=D$ are derivatives of rational functions; the
first term is a polynomial. They have zero residue at every finite point.
The remaining residue is $U_\alpha^{(D)}\Res(a_\alpha)$, nonzero at a
suitable point. Thus \eqref{eq:residue-coefficient} is not the zero
function. Clearing denominators gives a polynomial relation with
nonzero highest coefficient at index $\alpha-1$.
Descent eventually gives a nonzero polynomial times the nonzero
constant $v_0$ equal to zero, a contradiction.
\end{proof}

In \eqref{eq:lifted-system}, the link $-2/w$ has residue $-2$ at $0$,
and $4/(1-w^2)$ has residue $-2$ at $1$.
Consider a nonzero constant combination of the solutions $Y_q$.
If $q_0$ is the largest index with nonzero coefficient, its first
nonzero component is at index $h-q_0$ and is a nonzero constant,
by \eqref{eq:Yq}. This component remains constant after the change
of variable and rescaling.
Lemma~\ref{lem:chain}, starting at that index, shows that the remainder
cannot vanish identically, because its highest coefficient
$\widetilde S_{b+h}=w^{2M-1}P_b(w^2)$ is nonzero.
Thus the $h+1$ high-order remainders are linearly independent at
each of $\ii$ and $-\ii$.

\subsection{The other local remainders}

We also need remainders at $0$, infinity and $\pm1$. These will come from
solutions whose first $h+1$ components vanish. The remaining equations
then integrate to powers of a logarithm. Their combined orders at $0$
and infinity will be independent of $M$, while one of them will inherit
additional zeros at $\pm1$ from the decay of $F_n$.

For $j=1,\ldots,b$, system \eqref{eq:lifted-system} has a solution
whose components below $h+j$ vanish and whose remaining components are
\[
 \widetilde v_{i,j}(w)
     =\frac{(-\log(w^2))^{i-h-j}}{(i-h-j)!}
     \qquad(i\geq h+j).
\]
Its remainder is
\begin{equation}\label{eq:homogeneous-remainders}
 \widetilde\rho_j(w)=w^{2M-1}
       \sum_{i=j}^bP_i(w^2)
            \frac{(-\log(w^2))^{i-j}}{(i-j)!}.
\end{equation}
These $b$ remainders are independent.
In a nonzero constant relation, the least index $j$ with nonzero
coefficient supplies the unique highest logarithmic degree, with a
nonzero coefficient proportional to $P_b$.
This is impossible because $\log(w^2)$ is transcendental over $\C(w)$.

At $0$, each remainder has order at least $2M-1$; at infinity its order
is at least $-2(M+n)+1$.

For the zeros at $\pm1$, consider
\begin{equation}\label{eq:Rn-zero}
 R_n(z):=\sum_{i=1}^b P_i(z)\frac{(-\log z)^{i-1}}{(i-1)!}
       =O((z-1)^{\omega n-1})
\end{equation}
near $1$, on the branch $\log1=0$.
The asserted order is another form of the vanishing of $F_n$ at infinity.
Indeed, the partial fractions give the Laplace-transform identity
\[
 F_n(t)=\int_0^\infty e^{-tu}
       \sum_{i=1}^b P_i(e^{-u})\frac{u^{i-1}}{(i-1)!}\dd u
 \qquad(t>0).
\]
If the analytic function multiplying $e^{-tu}$ is $G(u)$, the coefficient
of $t^{-k-1}$ is $G^{(k)}(0)$. Proposition~\ref{prop:rational} therefore
implies that $G$ vanishes to order at least $\omega n-1$. Since
$G(u)=R_n(e^{-u})$, this proves \eqref{eq:Rn-zero}.
This is the observation used in \cite[Sections~3.5 and~4.5]{F26}.

Now $\widetilde\rho_1(w)=w^{2M-1}R_n(w^2)$.
At each of $w=1$ and $w=-1$, choose the local branch for which
$\log(w^2)=0$ at that point. The map $w\mapsto w^2$ has nonzero
derivative there, so \eqref{eq:Rn-zero} gives order at least
$\omega n-1$ at each point.
These two local branches need not be restrictions of a single global
branch; Proposition~\ref{prop:multiplicity} allows the different local
solutions.

\subsection{Counting the orders}

Apply Proposition~\ref{prop:multiplicity} with
\[
 \Sigma=\{0,\infty,1,-1,\ii,-\ii\},\qquad m=4M
\]
to \eqref{eq:lifted-system} and \eqref{eq:lifted-polynomials}.
The system is regular singular: its matrix has only simple finite poles
and is $O(1/w)$ at infinity. Its solutions therefore have the required
Nilsson expansions.
The constructed families give the following bounds.
\begin{center}
\begin{tabular}{@{}lcc@{}}
\toprule
Point & Independent remainders & Order of each, at least\\
\midrule
Each of $\ii,-\ii$ & $h+1$ & $\kappa n$\\
$0$ & $b$ & $2M-1$\\
$\infty$ & $b$ & $-2(M+n)+1$\\
Each of $1,-1$ & $1$ & $\omega n-1$\\
\bottomrule
\end{tabular}
\end{center}

There is a common rational differential operator of order
$\mu\leq b+h+1$ annihilating all remainders.
Differentiating a remainder iterates the operation
$S\mapsto S'+S\matA$ on a row of length $b+h+1$; its first $b+h+2$
iterates are rationally dependent.
Since the $b$ remainders at infinity are independent, $\mu\geq b$.
The multiplicity estimate gives
\begin{align*}
 \bigl(2(h+1)\kappa-2b+2\omega\bigr)n-2
   &\leq(4rn+1)(\mu-b)+C\\
   &\leq(4rn+1)(h+1)+C.
\end{align*}
The constant depends only on the system and $\Sigma$.
The finitely many possibilities $1\leq b\leq a$ allow a uniform constant,
even when $b$ varies with $n$.
Dividing the last inequality by $2n$ would give
\[
 (h+1)(\kappa-2r)+\omega\leq b+o(1)\leq a+o(1).
\]
This contradicts \eqref{eq:rank-condition} for sufficiently large $n$
and proves Proposition~\ref{prop:rank}.

